\section{交换式的范畴与观测算子表示}
\label{app:categories}

本附录给出交换式的范畴与观测算子表示，并说明带标签近似证书的复合律。
有限确定性 $\Gamma_N$ 判据继续以任务保真、状态压缩和闭合误差为附加条件。
Markov category 的系统性框架见\citet{Fritz2020}。

\subsection{精确交织形成动力学态射}

有限集合与 Markov 核构成范畴 $\mathbf{FinStoch}$：对象是有限集合，
态射复合为核乘法，恒等态射为确定性恒等核。普通映射
$q:X\to Y$ 嵌入为确定性核 $Q_q$。

把一个随机动力系统记为对象 $(X,K)$。从 $(X,K)$ 到 $(Y,L)$ 的精确动力学
态射是满足
\[
KQ=QL
\]
的核 $Q:X\rightsquigarrow Y$。若
$R:(Y,L)\to(Z,M)$ 也是精确态射，则
\[
K(QR)=(KQ)R=(QL)R=Q(LR)=Q(RM)=(QR)M,
\]
故 $QR$ 仍是精确态射。恒等核显然满足交换式，所以这些对象与态射
确实形成一个普通范畴。正文的确定性任务保真粗粒化对应其中
$Q=Q_q$ 且任务三角形交换的一类态射；严格信息商的存在还需要任务保真与
状态压缩条件。

\subsection{带收缩标签的近似证书及其复合律}

附录~\ref{app:composition-time} 已证明近似动力学因子的复合界。把闭合误差
上界与通道收缩系数记为标签 $(\eps,\rho)$ 时，复合标签为
\[
(\eps_Q,\rho_Q)\star(\eps_R,\rho_R)
=(\rho_R\eps_Q+\eps_R,\rho_Q\rho_R).
\]
该运算满足结合律，恒等通道的标签为 $(0,1)$。因此，精确动力学态射的普通
范畴可以扩充为带误差与收缩标签的证书系统；标签复合显式记录误差传播。

若另有任务通道 $G_X:X\rightsquigarrow O$、
$G_Y:Y\rightsquigarrow O$ 且
$d_\infty(G_X,QG_Y)\le\tau$，再与下一层任务证书复合，任务读出误差也可由
三角不等式累计。分别记录动力学闭合误差与任务读出误差可以保留两类证书的
语义；合并为单一数值则给出较粗的复合标签。

\subsection{有限确定性商的观测算子对偶}

Koopman 在确定性 Hamilton 系统中用线性算子推进观测量
\citep{Koopman1931}。下文的 $P_K$ 是这一观点对有限随机核的直接线性扩展，
所需性质由矩阵计算自足给出。

核 $K:X\rightsquigarrow Y$ 诱导线性观测算子
\[
P_K:\mathbb R^Y\to\mathbb R^X,
\qquad
(P_Kf)(x)=\sum_yK(x,y)f(y),
\]
且 $P_{KL}=P_KP_L$。确定性 $q:X\to Y$ 对应拉回
$P_{Q_q}f=f\circ q$。交换式 $KQ_q=Q_qL$ 等价于
\[
P_K(q^*f)=q^*(P_Lf)
\qquad\forall f\in\mathbb R^Y,
\]
即 $q^*\mathbb R^Y$ 是 $P_K$-不变的、含常数且对逐点乘法封闭的观测代数。
附录~\ref{app:exact-theory}证明，在有限集合上这类代数恰与分区对应。

因此，“宏观商空间”“闭合观测代数”和“交换方块”是同一有限精确结构的
三种等价表示。任务保真由任务因子化保证，压缩由相对状态数度量，近似最优值
达到依赖候选集的有限性或紧性，系统族满足 $\Gamma_N\to0$ 则依赖状态比例与
闭合误差的跨规模联合控制。
